All experiments are executed under a two-stage orchestration model.

\textbf{Stage 1} (\texttt{scripts/run\_full\_experiment.py}) performs generation and
scoring for each pair, mode, and seed.

\textbf{Stage 2} (\texttt{scripts/run\_post\_pipeline.py}) performs aggregation,
metric computation, statistical testing, figure generation, and manuscript assembly.

Fairness controls:
\begin{itemize}
\item identical seed set across modes;
\item identical molecule budget per run;
\item identical hit thresholds per pair from \texttt{configs/metrics.yaml};
\item identical score convention (lower is better);
\item run-level metadata and timestamped logs for traceability.
\end{itemize}

Any interrupted run is resumed only if required artifacts are missing; otherwise it is skipped.

\subsection*{Key comparisons}

\begin{itemize}
\item \texttt{double} vs \texttt{single\_a} cross-evaluated on target B:
  does \texttt{double} lose on one target for dual-activity?
\item \texttt{double} vs $\text{top}(\texttt{single\_a}) \cap \text{top}(\texttt{single\_b})$:
  does \texttt{double} find molecules that post-hoc intersection misses?
\item Homologous pairs (AChE\_BChE, CDK2\_ERK2) vs heterogeneous (BACE1\_GSK3B):
  how much does pocket similarity facilitate dual generation?
\item UniDock GPU ranking vs Boltz-2 ranking: are two independent methods in agreement?
\end{itemize}
